\slajdid{10-s052}
\begin{frame}[label=10-s052]{Inverzni režim rada tranzistora}
\gusttekst\setlength{\abovedisplayskip}{3pt}\setlength{\belowdisplayskip}{3pt}\setlength{\abovedisplayshortskip}{2pt}\setlength{\belowdisplayshortskip}{2pt}\begin{columns}[c,onlytextwidth]\column{.34\textwidth}\centering\begin{tikzpicture}[x=.65cm,y=.55cm,font=\sitneoznake,>=Latex]\ctikzset{bipoles/length=.55cm}\draw(0,-.5)--(0,.5);\draw(-2.3,0)node[ocirc]{}node[left]{$V_I$}to[R,l=$R_B$](-.6,0)--(0,0);\draw[->](0,.32)--(.7,.7);\draw(.7,.7)--(.7,1.3)to[R,l_=$R_C$](.7,3);\draw(.4,3)--(1,3);\node[above]at(.7,3){$V_{CC}$};\draw(.7,1.3)--(2,1.3)node[ocirc]{}node[right]{$V_o$};\draw(0,-.32)--(.7,-.7)--(.7,-1.5)node[ground]{};\node[right]at(.8,.6){C novi};\node[right]at(.8,-.65){E novi};\draw[->](-.7,.35)--(-.15,.35)node[midway,above]{$I_B$};\draw[->](.95,1.25)--(.95,.8)node[midway,right]{$I_C$};\draw[->](.95,-.95)--(.95,-1.45)node[midway,right]{$I_E$};\end{tikzpicture}
\par\medskip\begin{tikzpicture}[x=.65cm,y=.55cm,font=\sitneoznake,>=Latex]\draw[->](-.8,0)--(3.4,0)node[below]{$V_{CE}$};\draw[->](.7,-.5)--(.7,2.9)node[left]{$I_C$};\draw(.65,-.55)..controls(.75,.1)and(.88,1.7)..(1.05,2.1)..controls(1.2,2.5)and(2,2.5)..(2.85,2.65)node[right]{$V_{BE}$};\draw(.7,0)circle(.25);\draw[->](.95,.15)--(4.2,1.05);\node[above]at(2.6,.5){uvećano};\begin{scope}[shift={(4.45,.4)}]\draw(-.6,.8)--(1,.8);\draw(0,-.6)--(0,1.65);\draw(-.3,-.6)--(.45,1.45);\end{scope}\end{tikzpicture}
\column{.64\textwidth}
Emitor i kolektor zamene mesta. Tranzistor isto kao i u direktnom režimu može biti u inverznom aktivnom režimu kada važi da je bazno kolektorski spoj direktno polarizovan a bazno emiterski spoj inverzno polarizovan a veza struja je
\[I_{Cn}=\beta_R I_B\quad I_{En}=(\beta_R+1)I_B\quad0.1<\beta_R<5\]
gde je $\beta_R$ strujno pojačanje u inverznom aktivnom režimu. Inverzno zasićenje, uslov je
\[\beta_R I_{Bn}>I_{Cn}\qquad V_{CE}=V_{CESI}>0\]
gde je $V_{CESI}$ napon inverznog zasićenja. Obratite pažnju da su indeksi za „novi“ kolektor i emiter. Sa pravim indeksima
\[I_E=-\beta_R I_B\quad I_C=-(\beta_R+1)I_B\quad0.1<\beta_R<5\]
\[\beta_R I_B>|I_E|\qquad V_{CE}=V_{CESI}<0\]
\end{columns}
\end{frame}
